\chapter{The Machine}\label{cha:machine}

The BMC-K4510 is a fantasy computer: a machine that never existed, built the
way 1985 might have built it with no budget committee in the room. It is not
an emulation of any real computer. Its CPU, video chip, sound, and operating
software are its own; the parts it borrows --- a 6502-family instruction set,
the SID sound chip --- it borrows openly and then outgrows.

\section{What is in the box}
\begin{itemize}[leftmargin=1.4em]
\item \textbf{CPU: the 45GS10} --- the MEGA65's 45GS02 instruction set (a
  4510 with the Q pseudo-register and 32-bit flat addressing) plus this
  machine's own memory management: bank registers, a far-call gate, and RAM
  under the ROM. It runs at 40.5\,MHz.
\item \textbf{Memory: 256\,MB}, flat, 28-bit. The CPU sees 64\,KB at a time;
  everything else is one instruction away.
\item \textbf{Video: VICKe} --- 640$\times$480, 256 colours from a 24-bit
  palette, four layers, 128 sprites, a blitter that draws lines and filled
  triangles, and a display-list coprocessor named SHEILA.
\item \textbf{Sound: four SID chips} (reSID 6581) and a floating-point MATH
  unit the BASIC leans on.
\item \textbf{A system ROM}: a shell with directories, a Wozmon-style
  monitor, and EhBASIC~2.22 extended with graphics and sound --- with three
  more languages one word away: BBC BASIC on the Tube, Forth, and CP/M.
\end{itemize}

\section{Getting one}
On the \textbf{desktop} (Linux), three lines build the whole machine:
\begin{type}
git clone https://github.com/mlongval/bmc-k4510
cd bmc-k4510 && ./setup.sh
./k4510
\end{type}
\verb|setup.sh| installs the compilers and SDL2, builds everything, and
runs the machine's ten test suites.

On the \textbf{Raspberry Pi 3B+}, download the SD-card zip from the
project's releases page and put it on a card --- either unzip it onto the
card's root yourself, or let the script do it:
\begin{type}
./install-sd.sh bmc-k4510-pi3.zip /media/you/CARD
\end{type}

\begin{aside}
\textbf{Use an SDHC card --- 4 to 32\,GB, specifically.} SDHC ships
FAT32 from the factory and just works. Older plain SD cards (2\,GB and
under) do \emph{not} boot --- we tested. Bigger SDXC cards (64\,GB+)
ship exFAT and will not boot until reformatted;
\verb|./install-sd.sh zip /dev/sdX --format| does that (and erases the
card). The machine needs about 5\,MB, so the smallest SDHC card sold is
plenty.
\end{aside}

\section{First boot}
Switch it on (on the desktop: \verb|./k4510|; on the
Raspberry Pi 3B+: insert the card and power up). The machine greets you in yellow on blue:

\screen{shots/boot.png}{The boot screen: the hourglass, the name, the speed,
the memory. Everything else is one \texttt{INFO} away.}

The \verb|/]| at the bottom is the shell prompt --- the \verb|/| is the
directory you are in. Type \verb|HELP| for the commands, \verb|INFO| for the
machine's full self-description, and \verb|DIR| to see your files.

\begin{aside}
\textbf{Keys on the desktop:} Escape is RUN/STOP (stops a BASIC program),
Ctrl-C is STOP, Shift+Escape quits the emulator. Reset is a chord, the
way it was Commodore+Restore on the C64: Super+PageUp (the Windows or
Command key and PageUp) --- one key cannot do it by accident. F7 opens
the menu.
\end{aside}

\section{The F7 menu}
F7 (unshifted; Shift+F7 still reaches BBC BASIC) freezes the machine and
opens a framed window over the dimmed screen, in the manner of the
C64~Ultimate's: the picture behind it stays exactly as it was, sound
stops, and closing the menu resumes where it froze. The menu draws with
its own font and never touches the machine's memory, so it opens even
when a program has wrecked the screen. Cursor keys move, Enter selects,
Left/Right step a value, Escape backs out a level, F7 or Escape at the
top closes.

\screen{shots/menu.png}{The F7 menu's Video page with the screen-font
popup, over the frozen, dimmed boot screen.}

\begin{description}
\item[Video] border width and colour around the picture; the screen font
  (\emph{kernel8}, \emph{unscii}, \emph{open-roms}, \emph{PXLfont} ---
  a live swap of the chargen at \texttt{\$010000}, the two PETSCII fonts
  rearranged into ASCII order); full screen.
\item[Audio] volume.
\item[Input] the reset chord (Super, Ctrl or Alt + PageUp, or
  Ctrl+Alt+Del); which key opens the menu (F7, F8, F11 or Pause).
\item[Machine] \emph{save state} and \emph{load state}, four slots
  each (the whole machine --- CPU, every used page of the 256~MB, the
  banks and MAP, VICKe, the devices, JIM --- to \texttt{k4510-slotN.k4s}
  beside the settings file; the row shows when the slot was written;
  the Tube co-processor is not in the file and is stopped by a load);
  reset; power cycle (memory cleared, ROM reloaded); stop the Tube
  co-processor; quit.
\item[Info] version, ROM, files, host.
\end{description}

Leaving the menu writes the settings to \texttt{k4510.cfg} beside
\texttt{fs/} (on the Pi, \texttt{SD:/k4510/k4510.cfg}) --- a plain
\texttt{key = value} file you may edit; comments and keys it does not
know are kept. Adding a setting is one row in \texttt{core/ui/settings.c},
one row in the menu tree in \texttt{core/ui/menu.c}, and the place that
reads it; the same code runs on the Pi, where the C64 keyboard's F7 opens
it too.
